
\FigureLayoutDeclare{img/2020/shanghai/q10_fig1.png}{6.0cm}{8bp 6bp 6bp 8bp}
\FigureLayoutDeclare{img/2020/shanghai/q12_fig1.png}{8.0cm}{9bp 6bp 6bp 7bp}

\chapter{2020年上海卷（秋）}

\section{填空题}

\begin{problem}
已知集合 \(A=\{1,2,4\}\)，集合 \(B=\{2,4,5\}\)，则 \(A\cap B=\)\fillinblank{}.
\end{problem}

\begin{answer}
\(\{2,4\}\)
\end{answer}

\begin{solution}
因为 \(A=\{1,2,4\}\)，\(B=\{2,4,5\}\)，
\(A\cap B=\{2,4\}\).
故答案为 \(\{2,4\}\).
\end{solution}

\begin{problem}
计算：\(\lim\limits_{n\to\infty}\frac{n+1}{3n-1}=\)\fillinblank{}.
\end{problem}

\begin{answer}
\(\frac13\)
\end{answer}

\begin{solution}
因为
\(\lim\limits_{n\to\infty}\frac{n+1}{3n-1} =\lim\limits_{n\to\infty}\frac{1+\frac1n}{3-\frac1n} =\frac13\)，
故答案为 \(\frac13\).
\end{solution}

\begin{problem}
已知复数 \(z=1-2\i\)（\(\i\) 为虚数单位），则 \(|z|=\)\fillinblank{}.
\end{problem}

\begin{answer}
\(\sqrt5\)
\end{answer}

\begin{solution}
由 \(z=1-2\i\)，得
\(|z|=\sqrt{1^2+(-2)^2}=\sqrt5\).
故答案为 \(\sqrt5\).
\end{solution}

\begin{problem}
已知函数 \(f(x)=x^3\)，\(f^{-1}(x)\) 是 \(f(x)\) 的反函数，则 \(f^{-1}(x)=\)\fillinblank{}.
\end{problem}

\begin{answer}
\(\sqrt[3]{x}\)
\end{answer}

\begin{solution}
由 \(y=f(x)=x^3\)，得 \(x=\sqrt[3]{y}\).
把 \(x\) 与 \(y\) 互换，可得 \(f(x)=x^3\) 的反函数为
\(f^{-1}(x)=\sqrt[3]{x}\).
故答案为 \(\sqrt[3]{x}\).
\end{solution}

\begin{problem}
已知 \(x\)，\(y\) 满足 \(x+y-2\ge0\)，\(x+2y-3\le0\)，\(y\ge0\)，
则 \(z=y-2x\) 的最大值为\fillinblank{}.

\bitmapfigure[max width=6.0cm]{img/2020/shanghai/q5_fig1.png}
\end{problem}

\begin{answer}
\(-1\)
\end{answer}

\begin{solution}
由约束条件作出可行域如图阴影部分.
化目标函数 \(z=y-2x\) 为 \(y=2x+z\).
由图可知，当直线 \(y=2x+z\) 过点 \(A\) 时，直线在 \(y\) 轴上的截距最大.
联立
\examdisplaycases{}{
x+y-2=0,\\
x+2y-3=0,
}
解得 \(x=1\)，\(y=1\)，即 \(A(1,1)\).
所以 \(z\) 有最大值为 \(1-2\times1=-1\).
故答案为 \(-1\).
\end{solution}

\begin{problem}
已知行列式 \(\left|\begin{matrix}
1 & a & b\\
2 & c & d\\
3 & 0 & 0
\end{matrix}\right|=6\)，则 \(\left|\begin{matrix}
a & b\\
c & d
\end{matrix}\right|\)=\fillinblank{}.
\end{problem}

\begin{answer}
\(2\)
\end{answer}

\begin{solution}
由行列式 \(\left|\begin{matrix}
1 & a & b\\
2 & c & d\\
3 & 0 & 0
\end{matrix}\right|=6\)，可得 \(3\left|\begin{matrix}a & b\\ c & d\end{matrix}\right|=6\)，解得 \(\left|\begin{matrix}
a & b\\
c & d
\end{matrix}\right|=2\).
故答案为 \(2\).
\end{solution}

\begin{problem}
已知有四个数 \(1\)，\(2\)，\(a\)，\(b\)，这四个数的中位数是 \(3\)，平均数是 \(4\)，则 \(ab=\)\fillinblank{}.
\end{problem}

\begin{answer}
\(36\)
\end{answer}

\begin{solution}
因为四个数的平均数为 \(4\)，所以
\(a+b=4\times4-1-2=13\).
因为中位数是 \(3\)，所以
\(\frac{2+a}{2}=3\)，
解得 \(a=4\)，代入上式得 \(b=9\).
所以
\(ab=36\).
故答案为 \(36\).
\end{solution}

\begin{problem}
已知数列 \(\{a_n\}\) 是公差不为零的等差数列，且 \(a_1+a_{10}=a_9\)，则
\(\frac{a_1+a_2+\cdots+a_9}{a_{10}}\)=\fillinblank{}.
\end{problem}

\begin{answer}
\(\frac{27}{8}\)
\end{answer}

\begin{solution}
根据题意，等差数列 \(\{a_n\}\) 满足 \(a_1+a_{10}=a_9\)，即
\(a_1+a_1+9d=a_1+8d\)，
变形可得
\(a_1=-d\).
所以
\(\frac{a_1+a_2+\cdots+a_9}{a_{10}} =\frac{9a_1+36d}{a_1+9d} =\frac{9a_1+36d}{-d+9d} =\frac{27}{8}\).
故答案为 \(\frac{27}{8}\).
\end{solution}

\begin{problem}
从 \(6\) 个人挑选 \(4\) 个人去值班，每人值班一天，第一天安排 \(1\) 个人，第二天安排 \(1\) 个人，第三天安排 \(2\) 个人，则共有\fillinblank{} 种安排情况.
\end{problem}

\begin{answer}
\(180\)
\end{answer}

\begin{solution}
根据题意，可得排法共有
\(\binom{6}{1}\binom{5}{1}\binom{4}{2}=180\)
种.
故答案为 \(180\).
\end{solution}

\begin{problem}
已知椭圆 \(C:\frac{x^2}{4}+\frac{y^2}{3}=1\) 的右焦点为 \(F\)，直线 \(l\) 经过椭圆右焦点 \(F\)，交椭圆 \(C\) 于 \(P\)，\(Q\) 两点（点 \(P\) 在第二象限），若点 \(Q\) 关于 \(x\) 轴对称点为 \(Q'\)，且满足 \(PQ\perp FQ'\)，求直线 \(l\) 的方程是\fillinblank{}.

\bitmapfigure[max width=6.0cm]{img/2020/shanghai/q10_fig1.png}
\end{problem}

\begin{answer}
\(x+y-1=0\)
\end{answer}

\begin{solution}
椭圆 \(C:\frac{x^2}{4}+\frac{y^2}{3}=1\) 的右焦点为 \(F(1,0)\).
直线 \(l\) 经过椭圆右焦点 \(F\)，交椭圆 \(C\) 于 \(P\)，\(Q\) 两点（点 \(P\) 在第二象限）.
若点 \(Q\) 关于 \(x\) 轴对称点为 \(Q'\)，且满足 \(PQ\perp FQ'\)，
可知直线 \(l\) 的斜率为 \(-1\)，所以直线 \(l\) 的方程是
\(y=-(x-1)\)，
即
\(x+y-1=0\).
故答案为 \(x+y-1=0\).
\end{solution}

\begin{problem}
设 \(a\in\R\)，若存在定义域为 \(\R\) 的函数 \(f(x)\) 同时满足下列两个条件：

\begin{enumerate}
\item 对任意的 \(x_0\in\R\)，\(f(x_0)\) 的值为 \(x_0\) 或 \(x_0^2\)；
\item 关于 \(x\) 的方程 \(f(x)=a\) 无实数解，
\end{enumerate}

则 \(a\) 的取值范围是\fillinblank{}.
\end{problem}

\begin{answer}
\((-\infty,0)\cup(0,1)\cup(1,+\infty)\)
\end{answer}

\begin{solution}
根据条件（1）可得 \(f(0)=0\) 或 \(f(1)=1\).
又因为关于 \(x\) 的方程 \(f(x)=a\) 无实数解，所以 \(a\neq0\) 或 \(1\).
故
\(a\in(-\infty,0)\cup(0,1)\cup(1,+\infty)\).
故答案为 \((-\infty,0)\cup(0,1)\cup(1,+\infty)\).
\end{solution}

\begin{problem}
已知 \(\bs{a}_1\)，\(\bs{a}_2\)，\(\bs{b}_1\)，\(\bs{b}_2\)，\(\ldots\)，\(\bs{b}_k\)（\(k\in\N^*\)）是平面内两两互不相等的向量，满足 \(|\bs{a}_1-\bs{a}_2|=1\)，且 \(|\bs{a}_i-\bs{b}_j|\in\{1,2\}\)（其中 \(i=1\)，\(2\)，\(j=1\)，\(2\)，\(\ldots\)，\(k\)），则 \(k\) 的最大值是\fillinblank{}.

\bitmapfigure[max width=6.0cm]{img/2020/shanghai/q12_fig1.png}
\end{problem}

\begin{answer}
\(6\)
\end{answer}

\begin{solution}
如图，设 \(\overrightarrow{OA_1}=\bs{a}_1\)，\(\overrightarrow{OA_2}=\bs{a}_2\).
由 \(|\bs{a}_1-\bs{a}_2|=1\)，且 \(|\bs{a}_i-\bs{b}_j|\in\{1,2\}\)，分别以 \(A_1\)，\(A_2\) 为圆心，以 \(1\) 和 \(2\) 为半径画圆，其中任意两圆的公共点共有 \(6\) 个.
故满足条件的 \(k\) 的最大值为 \(6\).
故答案为 \(6\).
\end{solution}
\section{单选题}

\begin{problem}
下列不等式恒成立的是（\quad）

\choices
    {\(a^2+b^2\le2ab\)}
    {\(a^2+b^2\ge-2ab\)}
    {\(a+b\ge2\sqrt{|ab|}\)}
    {\(a^2+b^2\le-2ab\)}
\end{problem}

\begin{answer}
B
\end{answer}

\begin{solution}
A. 当 \(a<0\)，\(b>0\) 时，不等式 \(a^2+b^2\le2ab\) 不成立，故 A 错误；

B. 因为 \((a+b)^2\ge0\)，所以 \(a^2+b^2+2ab\ge0\)，即 \(a^2+b^2\ge-2ab\)，故 B 正确；

C. 当 \(a<0\)，\(b<0\) 时，不等式 \(a+b\ge2\sqrt{|ab|}\) 不成立，故 C 错误；

D. 当 \(a>0\)，\(b>0\) 时，不等式 \(a^2+b^2\le-2ab\) 不成立，故 D 错误.

故选 B.
\end{solution}

\begin{problem}
已知直线方程 \(3x+4y+1=0\) 的一个参数方程可以是（\quad）

\choices
    {\(x=1+3t\)，\(y=-1-4t\)（\(t\) 为参数）}
    {\(x=1-4t\)，\(y=-1+3t\)（\(t\) 为参数）}
    {\(x=1-3t\)，\(y=-1+4t\)（\(t\) 为参数）}
    {\(x=1+4t\)，\(y=1-3t\)（\(t\) 为参数）}
\end{problem}

\begin{answer}
B
\end{answer}

\begin{solution}
B 选项的普通方程为
\(\frac{x-1}{-4}=\frac{y+1}{3}\)，
即
\(3x+4y+1=0\)，
正确.

A 选项的普通方程为 \(4x+3y-1=0\)，不正确；
C 选项的普通方程为 \(4x+3y-1=0\)，不正确；
D 选项的普通方程为 \(3x+4y-7=0\)，不正确.

故选 B.
\end{solution}

\begin{problem}
在棱长为 \(10\) 的正方体 \(ABCD-A_1B_1C_1D_1\) 中，\(P\) 为左侧面 \(ADD_1A_1\) 上一点，已知点 \(P\) 到 \(A_1D_1\) 的距离为 \(3\)，\(P\) 到 \(AA_1\) 的距离为 \(2\)，则过点 \(P\) 且与 \(A_1C\) 平行的直线交正方体于 \(P\)，\(Q\) 两点，则 \(Q\) 点所在的平面是（\quad）

\bitmapfigure[max width=6.0cm]{img/2020/shanghai/q15_fig1.png}

\choices
    {\(AA_1B_1B\)}
    {\(BB_1C_1C\)}
    {\(CC_1D_1D\)}
    {\(ABCD\)}
\end{problem}

\begin{answer}
D
\end{answer}

\begin{solution}
如图，
\solutionfigure{\bitmapfigure[max width=3.8cm]{img/2020/shanghai/q15_solution_fig1.png}}
由点 \(P\) 到 \(A_1D_1\) 的距离为 \(3\)，\(P\) 到 \(AA_1\) 的距离为 \(2\)，
可得 \(P\) 在 \(\triangle AA_1D\) 内，过 \(P\) 作 \(EF\parallel A_1D\)，且 \(EF\cap AA_1\) 于 \(E\)，\(EF\cap AD\) 于 \(F\).
在平面 \(ABCD\) 中，过 \(F\) 作 \(FG\parallel CD\)，交 \(BC\) 于 \(G\)，则平面 \(EFG\parallel\) 平面 \(A_1DC\).
连接 \(AC\)，交 \(FG\) 于 \(M\)，连接 \(EM\).
因为平面 \(EFG\parallel\) 平面 \(A_1DC\)，所以平面 \(A_1AC\cap\) 平面 \(A_1DC=A_1C\)，平面 \(A_1AC\cap\) 平面 \(EFM=EM\)，
所以 \(EM\parallel A_1C\).
在 \(\triangle EFM\) 中，过 \(P\) 作 \(PQ\parallel EM\)，且 \(PQ\cap FM\) 于 \(Q\)，则 \(PQ\parallel A_1C\).
因为线段 \(FM\) 在四边形 \(ABCD\) 内，\(Q\) 在线段 \(FM\) 上，所以 \(Q\) 在四边形 \(ABCD\) 内.
所以 \(Q\) 点所在的平面是平面 \(ABCD\).

故选 D.
\end{solution}

\begin{problem}
命题 \(p\)：存在 \(a\in\R\) 且 \(a\neq0\)，对于任意的 \(x\in\R\)，使得 \(f(x+a)<f(x)+f(a)\)；

命题 \(q_1\)：\(f(x)\) 单调递减且 \(f(x)>0\) 恒成立；

命题 \(q_2\)：\(f(x)\) 单调递增，存在 \(x_0<0\) 使得 \(f(x_0)=0\)，

则下列说法正确的是（\quad）

\choices
    {只有 \(q_1\) 是 \(p\) 的充分条件}
    {只有 \(q_2\) 是 \(p\) 的充分条件}
    {\(q_1\)，\(q_2\) 都是 \(p\) 的充分条件}
    {\(q_1\)，\(q_2\) 都不是 \(p\) 的充分条件}
\end{problem}

\begin{answer}
C
\end{answer}

\begin{solution}
对于命题 \(q_1\)：当 \(f(x)\) 单调递减且 \(f(x)>0\) 恒成立时，
当 \(a>0\) 时，此时 \(x+a>x\)，又因为 \(f(x)\) 单调递减，所以 \(f(x+a)<f(x)\).
又因为 \(f(x)>0\) 恒成立，所以 \(f(x)<f(x)+f(a)\)，从而
\(f(x+a)<f(x)+f(a)\).
所以命题 \(q_1\Rightarrow p\).

对于命题 \(q_2\)：当 \(f(x)\) 单调递增，存在 \(x_0<0\) 使得 \(f(x_0)=0\) 时，
取 \(a=x_0<0\)，此时 \(x+a<x\)，又因为 \(f(x)\) 单调递增，所以
\(f(x+a)<f(x)\).
同时 \(f(a)=f(x_0)=0\)，故
\(f(x+a)<f(x)+f(a)\).
所以命题 \(q_2\Rightarrow p\).

所以 \(q_1\)，\(q_2\) 都是 \(p\) 的充分条件.
故选 C.
\end{solution}
\section{解答题}

\begin{problem}
已知 \(ABCD\) 是边长为 \(1\) 的正方形，正方形 \(ABCD\) 绕 \(AB\) 旋转形成一个圆柱.

\begin{enumerate}
\item 求该圆柱的表面积；
\item 正方形 \(ABCD\) 绕 \(AB\) 逆时针旋转 \(\frac{\pi}{2}\) 至 \(ABC_1D_1\)，求线段 \(CD_1\) 与平面 \(ABCD\) 所成的角.
\end{enumerate}

\bitmapfigure[max width=6.0cm]{img/2020/shanghai/q17_fig1.png}
\end{problem}

\begin{solution}
\begin{enumerate}
\item 该圆柱的表面由上下两个半径为 \(1\) 的圆面和一个长为 \(2\pi\)，宽为 \(1\) 的矩形组成，
\(S=2\times\pi\times1^2+2\pi\times1=4\pi\).
故该圆柱的表面积为 \(4\pi\).

\item 因为正方形 \(ABC_1D_1\)，所以 \(AD_1\perp AB\).
又 \(\angle DAD_1=\frac{\pi}{2}\)，所以 \(AD_1\perp AD\).
因为 \(AD\cap AB=A\)，且 \(AD\)，\(AB\subset\) 平面 \(ADB\)，所以 \(AD_1\perp\) 平面 \(ADB\)，即 \(D_1\) 在面 \(ADB\) 上的投影为 \(A\).
连接 \(CD_1\)，则 \(\angle D_1CA\) 即为线段 \(CD_1\) 与平面 \(ABCD\) 所成的角.
由题意可得
\(\cos\angle D_1CA=\frac{AC}{CD_1}=\frac{\sqrt2}{\sqrt3}=\frac{\sqrt6}{3}\).
所以线段 \(CD_1\) 与平面 \(ABCD\) 所成的角为
\(\arccos\frac{\sqrt6}{3}\).
\end{enumerate}
\end{solution}

\begin{problem}
已知函数 \(f(x)=\sin\omega x\)，\(\omega>0\).

\begin{enumerate}
\item \(f(x)\) 的周期是 \(4\pi\)，求 \(\omega\)，并求 \(f(x)=\frac12\) 的解集；
\item 已知 \(\omega=1\)，\(g(x)=f^2(x)+\sqrt3\,f(-x)f\!\left(\frac\pi2-x\right)\)，\(x\in\left[0,\frac\pi4\right]\)，求 \(g(x)\) 的值域.
\end{enumerate}
\end{problem}

\begin{solution}
\begin{enumerate}
\item 由于 \(f(x)\) 的周期是 \(4\pi\)，所以
\(\omega=\frac{2\pi}{4\pi}=\frac12\)，
所以
\(f(x)=\sin\frac{x}{2}\).
令
\(\sin\frac{x}{2}=\frac12\)，
故
\(\frac{x}{2}=\frac{\pi}{6}+2k\pi\)或\(\frac{x}{2}=\frac{5\pi}{6}+2k\pi\)，
整理得
\(x=4k\pi+\frac{\pi}{3}\)或\(x=4k\pi+\frac{5\pi}{3}\)，\(k\in\Z\).

\item 由于 \(\omega=1\)，所以 \(f(x)=\sin x\).
所以
\(g(x)=\sin^2x+\sqrt3\sin(-x)\sin\left(\frac\pi2-x\right)\)
\(=\frac{1-\cos2x}{2}-\frac{\sqrt3}{2}\sin2x\)
\(=\frac12-\frac12\left(\cos2x+\sqrt3\sin2x\right)\)
\(=\frac12-\sin\left(2x+\frac\pi6\right)\).
由于 \(x\in\left[0,\frac\pi4\right]\)，所以
\(\frac\pi6\le2x+\frac\pi6\le\frac{2\pi}{3}\).
于是
\(\frac12\le\sin\left(2x+\frac\pi6\right)\le1\)，
故
\(-1\le-\sin\left(2x+\frac\pi6\right)\le-\frac12\).
故
\(-\frac12\le g(x)\le0\).
所以函数 \(g(x)\) 的值域为 \(\left[-\frac12,0\right]\).
\end{enumerate}
\end{solution}

\begin{problem}
在研究某市交通情况时，道路密度是指该路段上一定时间内通过的车辆数除以时间，车辆密度是该路段一定时间内通过的车辆数除以该路段的长度，现定义交通流量 \(v=f(x)=\frac{q}{x}\)，其中 \(x\) 为道路密度，\(q\) 为车辆密度，且
\(f(x)=100-135\left(\frac13\right)^{\frac{80}{x}}\)（\(0<x<40\)），
\(f(x)=-k(x-40)+85\)（\(40\le x\le80\)）.

\begin{enumerate}
\item 若交通流量 \(v>95\)，求道路密度 \(x\) 的取值范围；
\item 已知道路密度 \(x=80\) 时，测得交通流量 \(v=50\)，求车辆密度 \(q\) 的最大值.
\end{enumerate}
\end{problem}

\begin{solution}
\begin{enumerate}
\item 按实际情况而言，交通流量 \(v\) 随着道路密度 \(x\) 的增大而减小，故 \(v=f(x)\) 是单调递减函数，所以 \(k>0\).
当 \(40\le x\le80\) 时，\(v\) 最大为 \(85\).
于是只需令
\(100-135\left(\frac13\right)^{\frac{80}{x}}>95\)，
解得
\(\left(\frac13\right)^{\frac{80}{x}}<\frac1{27}\)，
所以
\(\frac{80}{x}>3\)，
即
\(x<\frac{80}{3}\).
故道路密度 \(x\) 的取值范围为 \(\left(0,\frac{80}{3}\right)\).

\item 把 \(x=80\)，\(v=50\) 代入 \(v=f(x)=-k(x-40)+85\) 中，
\(50=-k\cdot40+85\)，
解得
\(k=\frac78\).
所以
\[
q=vx=
\begin{cases}
x\left(100-135\left(\frac13\right)^{\frac{80}{x}}\right), & 0<x<40,\\
\left(-\frac78(x-40)+85\right)x, & 40\le x\le80.
\end{cases}
\]
当 \(0<x<40\) 时，
\(v=100-135\left(\frac13\right)^{\frac{80}{x}}<100\)，
所以
\(q=vx<100\times40=4000\).

当 \(40\le x\le80\) 时，
\(q=-\frac78x^2+120x\)，
对称轴为
\(x=\frac{480}{7}\)，
此时 \(q\) 有最大值，为
\(\frac{28800}{7}\)，
且 \(\frac{28800}{7}>4000\).

综上所述，车辆密度 \(q\) 的最大值为
\(\frac{28800}{7}\).
\end{enumerate}
\end{solution}

\begin{problem}
已知双曲线 \(\Gamma_1:\frac{x^2}{4}-\frac{y^2}{b^2}=1\) 与圆 \(\Gamma_2:x^2+y^2=4+b^2\)（\(b>0\)）交于点 \(A(x_A,y_A)\)（第一象限），曲线 \(\Gamma\) 为 \(\Gamma_1\)，\(\Gamma_2\) 上取满足 \(x>|x_A|\) 的部分.

\begin{enumerate}
\item 若 \(x_A=\sqrt6\)，求 \(b\) 的值；
\item 当 \(b=\sqrt5\)，\(\Gamma_2\) 与 \(x\) 轴交点记作点 \(F_1\)，\(F_2\)，\(P\) 是曲线 \(\Gamma\) 上一点，且在第一象限，且 \(|PF_1|=8\)，求 \(\angle F_1PF_2\)；
\item 过点 \(D\left(0,\frac{b^2}{2}+2\right)\) 斜率为 \(-\frac{b}{2}\) 的直线 \(l\) 与曲线 \(\Gamma\) 只有两个交点，记为 \(M\)，\(N\)，用 \(b\) 表示 \(\overrightarrow{OM}\cdot\overrightarrow{ON}\)，并求 \(\overrightarrow{OM}\cdot\overrightarrow{ON}\) 的取值范围.
\end{enumerate}
\end{problem}

\begin{solution}
\begin{enumerate}
\item 由 \(x_A=\sqrt6\)，点 \(A\) 为曲线 \(\Gamma_1\) 与曲线 \(\Gamma_2\) 的交点，联立
\examdisplaycases{}{
\frac{x_A^2}{4}-\frac{y_A^2}{b^2}=1,\\
x_A^2+y_A^2=4+b^2,
}
解得
\(y_A=\sqrt2\)，\(b=2\).

\item 由题意可得 \(F_1\)，\(F_2\) 为曲线 \(\Gamma_1\) 的两个焦点.
由双曲线的定义可得
\(|PF_1|-|PF_2|=2a\)，
又 \(|PF_1|=8\)，\(2a=4\)，所以
\(|PF_2|=8-4=4\).
因为 \(b=\sqrt5\)，则
\(c=\sqrt{4+5}=3\)，
所以
\(|F_1F_2|=6\).
在 \(\triangle PF_1F_2\) 中，由余弦定理可得
\(\cos\angle F_1PF_2=\frac{|PF_1|^2+|PF_2|^2-|F_1F_2|^2}{2|PF_1|\cdot|PF_2|}=\frac{64+16-36}{2\times8\times4}=\frac{11}{16}\).
由 \(0<\angle F_1PF_2<\pi\)，可得
\(\angle F_1PF_2=\arccos\frac{11}{16}\).

\item 设直线 \(l:y=-\frac{b}{2}x+\frac{b^2}{2}+2\)，可得原点 \(O\) 到直线 \(l\) 的距离
\(d=\frac{\frac{b^2+4}{2}}{\sqrt{1+\frac{b^2}{4}}}=\sqrt{4+b^2}\)，
所以直线 \(l\) 是圆的切线，设切点为 \(M\).

所以 \(k_{OM}=\frac{2}{b}\)，并设 \(OM:y=\frac{2}{b}x\) 与圆 \(x^2+y^2=4+b^2\) 联立，可得
\(x^2+\frac{4}{b^2}x^2=4+b^2\)，
\(x=b\)，\(y=2\)，
即 \(M(b,2)\).
注意直线 \(l\) 与双曲线的斜率为负的渐近线平行，所以只有当 \(y_A>2\) 时，直线 \(l\) 才能与曲线 \(\Gamma\) 有两个交点.
由
\examdisplaycases{}{
\frac{x_A^2}{4}-\frac{y_A^2}{b^2}=1,\\
x_A^2+y_A^2=4+b^2
}
可得
\(y_A^2=\frac{b^4}{4+b^2}\).
所以有
\(4<\frac{b^4}{4+b^2}\)，
解得
\(b^2>2+2\sqrt5\)
或 \(b^2<2-2\sqrt5\)（舍去）.
因为 \(\overrightarrow{OM}\) 为 \(\overrightarrow{ON}\) 在 \(\overrightarrow{OM}\) 上的投影，可得
\(\overrightarrow{OM}\cdot\overrightarrow{ON}=4+b^2\).
所以
\(\overrightarrow{OM}\cdot\overrightarrow{ON}>6+2\sqrt5\)，
则
\(\overrightarrow{OM}\cdot\overrightarrow{ON}\in(6+2\sqrt5,+\infty)\).
\end{enumerate}
\end{solution}

\begin{problem}
已知数列 \(\{a_n\}\) 为有限数列，满足 \(|a_1-a_2|\le|a_1-a_3|\le\cdots\le|a_1-a_m|\)，则称 \(\{a_n\}\) 满足性质 \(P\).

\begin{enumerate}
\item 判断数列 \(3\)，\(2\)，\(5\)，\(1\) 和 \(4\)，\(3\)，\(2\)，\(5\)，\(1\) 是否具有性质 \(P\)，请说明理由；
\item 若 \(a_1=1\)，公比为 \(q\) 的等比数列，项数为 \(10\)，具有性质 \(P\)，求 \(q\) 的取值范围；
\item 若 \(\{a_n\}\) 是 \(1\)，\(2\)，\(3\)，\(\ldots\)，\(m\) 的一个排列（\(m\ge4\)），\(\{b_n\}\) 符合 \(b_k=a_{k+1}\)（\(k=1\)，\(2\)，\(\ldots\)，\(m-1\)），\(\{a_n\}\)，\(\{b_n\}\) 都具有性质 \(P\)，求所有满足条件的数列 \(\{a_n\}\).
\end{enumerate}
\end{problem}

\begin{solution}
\begin{enumerate}
\item 对于数列 \(3\)，\(2\)，\(5\)，\(1\)，有 \(|2-3|=1\)，\(|5-3|=2\)，\(|1-3|=2\)，满足题意，该数列满足性质 \(P\).

对于第二个数列 \(4\)，\(3\)，\(2\)，\(5\)，\(1\)，有 \(|3-4|=1\)，\(|2-4|=2\)，\(|5-4|=1\)，不满足题意，该数列不满足性质 \(P\).

\item 由题意，
\(|a_1-a_1q^n|\ge|a_1-a_1q^{n-1}|\)，\(n\in\{2,3,\ldots,9\}\)，
可得
\(|q^n-1|\ge|q^{n-1}-1|\).
两边平方可得
\(q^{2n}-2q^n+1\ge q^{2n-2}-2q^{n-1}+1\).
整理可得
\((q-1)q^{n-1}\big[q^{n-1}(q+1)-2\big]\ge0\).
当 \(q\ge1\) 时，得 \(q^{n-1}(q+1)-2\ge0\)，此时关于 \(n\) 恒成立，所以等价于 \(n=2\) 时 \(q(q+1)-2\ge0\).
所以
\((q+2)(q-1)\ge0\)，
所以 \(q\le-2\) 或 \(q\ge1\)，所以取 \(q\ge1\).

当 \(0<q\le1\) 时，得 \(q^{n-1}(q+1)-2\le0\)，此时关于 \(n\) 恒成立，所以等价于 \(n=2\) 时 \(q(q+1)-2\le0\).
所以
\((q+2)(q-1)\le0\)，
所以 \(-2\le q\le1\)，所以取 \(0<q\le1\).

当 \(-1\le q<0\) 时矛盾，舍去.

当 \(q<-1\) 时，得
\(q^{n-1}\big[q^{n-1}(q+1)-2\big]\le0\).
当 \(n\) 为奇数时，得 \(q^{n-1}(q+1)-2\le0\)，恒成立；当 \(n\) 为偶数时，\(q^{n-1}(q+1)-2\ge0\)，恒成立. 故等价于 \(n=2\) 时 \(q(q+1)-2\ge0\).
所以
\((q+2)(q-1)\ge0\)，
所以 \(q\le-2\) 或 \(q\ge1\)，所以取 \(q\le-2\).

当 \(q=0\) 时，数列为 \(1,0,0,\ldots\)，所有距离均为 \(1\)，也满足性质 P. 综上
\(q\in(-\infty,-2]\cup[0,+\infty)\).

\item 设 \(a_1=p\)，\(p\in\{3,4,\ldots,m-3,m-2\}\).
因为 \(a_1=p\)，\(a_2\) 可以取 \(p-1\) 或 \(p+1\)，\(a_3\) 可以取 \(p-2\) 或 \(p+2\).
如果 \(a_2\) 或 \(a_3\) 取了 \(p-3\) 或 \(p+3\)，将使 \(\{a_n\}\) 不满足性质 \(P\)；所以 \(\{a_n\}\) 的前 \(5\) 项有以下组合：

\begin{enumerate}
\item \(a_1=p\)，\(a_2=p-1\)，\(a_3=p+1\)，\(a_4=p-2\)，\(a_5=p+2\)；

\item \(a_1=p\)，\(a_2=p-1\)，\(a_3=p+1\)，\(a_4=p+2\)，\(a_5=p-2\)；

\item \(a_1=p\)，\(a_2=p+1\)，\(a_3=p-1\)，\(a_4=p-2\)，\(a_5=p+2\)；

\item \(a_1=p\)，\(a_2=p+1\)，\(a_3=p-1\)，\(a_4=p+2\)，\(a_5=p-2\).
\end{enumerate}

\begin{enumerate}
\item 对于 \(b_1=p-1\)，\(|b_2-b_1|=2\)，\(|b_3-b_1|=1\)，与 \(\{b_n\}\) 满足性质 \(P\) 矛盾，舍去；

\item 对于 \(b_1=p-1\)，\(|b_2-b_1|=2\)，\(|b_3-b_1|=3\)，\(|b_4-b_1|=1\)，与 \(\{b_n\}\) 满足性质 \(P\) 矛盾，舍去；

\item 对于 \(b_1=p+1\)，\(|b_2-b_1|=2\)，\(|b_3-b_1|=3\)，\(|b_4-b_1|=1\)，与 \(\{b_n\}\) 满足性质 \(P\) 矛盾，舍去；

\item 对于 \(b_1=p+1\)，\(|b_2-b_1|=2\)，\(|b_3-b_1|=1\)，与 \(\{b_n\}\) 满足性质 \(P\) 矛盾，舍去.
\end{enumerate}

所以 \(p\in\{3,4,\ldots,m-3,m-2\}\) 时，均不能同时使 \(\{a_n\}\)，\(\{b_n\}\) 都具有性质 \(P\).

当 \(p=1\) 时，有数列 \(\{a_n\}:1\)，\(2\)，\(3\)，\(\ldots\)，\(m-1\)，\(m\) 满足题意.

当 \(p=m\) 时，有数列 \(\{a_n\}:m\)，\(m-1\)，\(\ldots\)，\(3\)，\(2\)，\(1\) 满足题意.

当 \(p=2\) 时，有数列 \(\{a_n\}:2\)，\(1\)，\(3\)，\(\ldots\)，\(m-1\)，\(m\) 满足题意.

当 \(p=m-1\) 时，有数列 \(\{a_n\}:m-1\)，\(m\)，\(m-2\)，\(m-3\)，\(\ldots\)，\(3\)，\(2\)，\(1\) 满足题意.

所以满足题意的数列 \(\{a_n\}\) 只有以上四种.
\end{enumerate}
\end{solution}
